% ==============================================================================
% SECTION 4: DISCUSSION
% ==============================================================================

\section{Discusión}\label{sec:discussion}

\subsection{Hallazgos principales}\label{subsec:main-findings}

Los hallazgos de esta revisión subrayan que los sistemas agénticos no son solo
una mejora incremental, sino un cambio de paradigma en la bioinformática
computacional. Se identificó que la autocorrección iterativa es el componente
vital que permite la automatización de tareas que anteriormente requerían
supervisión humana constante, como el modelado PopPK o ciertos análisis de
\textit{single-cell omics}, aunque la validación experta continúa siendo
necesaria en etapas críticas.

Un hallazgo fundamental es el valor de la orquestación jerárquica, donde un
agente supervisor coordina a especialistas en herramientas y distribuye tareas
analíticas de forma más estructurada. No obstante, los resultados también
revelan que la calidad de los prompts de sistema y la estructura de las
herramientas disponibles, es decir, herramientas agénticas o legibles por
máquina, son determinantes centrales del éxito del workflow.

Este patrón valida la necesidad de protocolos como el \textit{Model Context
Protocol} (MCP) para el futuro del área, ya que la interoperabilidad entre
agentes, herramientas bioinformáticas y entornos de ejecución científica emerge
como una condición central para escalar estos sistemas de forma reproducible.

\subsection{Limitaciones}\label{subsec:discussion-limitations}

A pesar del potencial disruptivo, la implementación de agentes en
bioinformática enfrenta barreras técnicas críticas. El comportamiento no
determinista de los modelos subyacentes introduce una variabilidad inaceptable
en flujos de trabajo que requieren precisión absoluta, como el modelado PopPK.
Asimismo, la propagación de errores en arquitecturas jerárquicas, donde una
alucinación inicial contamina las etapas subsiguientes, subraya la inmadurez de
los mecanismos de autocorrección actuales.

La literatura también destaca la ausencia de estándares de diseño, lo que genera
una fragmentación de herramientas que dificulta la interoperabilidad. Finalmente,
la reproducibilidad sigue siendo un desafío abierto; la naturaleza de ``caja
negra'' de los LLMs comerciales impide asegurar que un experimento automatizado
produzca resultados idénticos en el tiempo, lo que requiere una transición hacia
modelos locales y registros de observabilidad más estrictos.

\begin{itemize}
  \item \textbf{Comportamiento no determinista:} Los LLMs pueden generar
  respuestas distintas ante el mismo prompt, lo que en bioinformática es
  crítico. Un agente puede proponer un script de PyMOL funcional hoy y uno con
  errores de sintaxis mañana. Esto obliga a implementar capas de validación de
  esquemas, como propone \textit{ChatSpatial}, para forzar una estructura en la
  salida.

  \item \textbf{Propagación de errores:} En sistemas multi-agente, si el agente
  investigador extrae un dato erróneo de PubMed, el agente programador escribirá
  código basado en ese error y el agente analista generará conclusiones falsas.
  Sin un monitor de errores, como
  LangSmith~\citep{LangSmithObservability2026,LangSmithConcepts2026}, este
  error puede permanecer invisible hasta el final del flujo.

  \item \textbf{Falta de estándares de diseño:} Actualmente no existe un
  lenguaje universal para que los agentes interactúen con herramientas
  bioinformáticas. Cada desarrollador crea sus propias funciones de Python. Aquí
  surge la necesidad del \textit{Model Context Protocol}
  (MCP)~\citep{41729821} para estandarizar la invocación de herramientas
  bioinformáticas externas por parte de agentes.

  \item \textbf{Dificultad en la reproducibilidad:} Debido a la actualización
  constante de modelos comerciales, un workflow agéntico documentado en 2025
  podría no funcionar igual en 2026, lo que desafía los principios de la ciencia
  abierta y dificulta la revisión por pares.
\end{itemize}

\begin{table}[htbp]
  \centering
  \caption{Impacto de las principales limitaciones de los sistemas agénticos en
  bioinformática y estrategias de mitigación sugeridas.}
  \label{tab:discussion-limitations}
  \begin{tabularx}{\textwidth}{@{} l X X @{}}
    \toprule
    \textbf{Limitación} & \textbf{Impacto en bioinformática} &
    \textbf{Mitigación sugerida} \\
    \midrule
    No determinismo &
    Variabilidad en resultados de simulaciones. &
    Fijación de temperatura ($T=0$) y \textit{prompt engineering} robusto. \\
    Propagación &
    Alucinaciones científicas verosímiles. &
    Mecanismos de \textit{grounding} y validación cruzada entre agentes. \\
    Falta de estándares &
    Fragmentación de herramientas y silos. &
    Adopción de protocolos de interoperabilidad, como
    MCPmed~\citealp{41729821}. \\
    Reproducibilidad &
    Imposibilidad de replicar pipelines exactos. &
    Registro exhaustivo de versiones y trazas de ejecución, como
    LangSmith~\citealp{LangSmithObservability2026}. \\
    \bottomrule
  \end{tabularx}
\end{table}

\subsection{Vacíos en la literatura}\label{subsec:literature-gaps}

La revisión identifica vacíos críticos que limitan la adopción masiva de estas
tecnologías. En primer lugar, la falta de benchmarks estandarizados impide una
comparación justa entre las crecientes arquitecturas agénticas, relegando la
evaluación a éxitos anecdóticos en lugar de métricas universales.

En segundo lugar, existe una carencia de evaluaciones comparativas rigurosas que
analicen la relación entre el costo computacional y la ganancia en precisión
científica, un factor determinante para la sostenibilidad en el entorno
académico. Finalmente, la ausencia de estudios en entornos reales deja abierta
la capacidad de estos sistemas para manejar ambigüedad y requisitos de
seguridad sobre datos genómicos sensibles.

Cerrar estos vacíos será clave para pasar de prototipos experimentales a
sistemas de producción científica confiables.
